\documentclass[11pt,a4paper]{article}
\usepackage[italian]{babel}
\usepackage[margin=2.5cm]{geometry}
\usepackage{parskip}
\usepackage{amsmath, amssymb, amsthm}
\usepackage{booktabs}
\usepackage{array}
\usepackage{xcolor}
\usepackage{needspace}
\usepackage{enumitem}
\usepackage{listings}
\usepackage{tikz}
\usepackage[most]{tcolorbox}
\usepackage[hidelinks]{hyperref}
\usetikzlibrary{decorations.pathreplacing, shapes.geometric, arrows.meta, positioning}

% Niente vedove/orfane: i paragrafi non lasciano righe isolate a fine/inizio pagina
\widowpenalty=10000
\clubpenalty=10000

% Elenchi più compatti
\setlist{itemsep=2pt, parsep=0pt, topsep=4pt, partopsep=0pt}

% --- Riquadri con bordo nero, senza numero (si spezzano tra due pagine) ---
% Uso: \begin{definizione} ... \end{definizione}
%      \begin{teorema}[Nome] ... \end{teorema}  (il nome è facoltativo)
\tcbset{nota/.style={enhanced jigsaw, breakable, colback=white,
  colframe=black, boxrule=0.5pt, sharp corners}}
\NewTColorBox{definizione}{}{nota,
  before upper={\textbf{Definizione.}\ }}
\NewTColorBox{teorema}{o}{nota,
  before upper={\textbf{Teorema\IfValueT{#1}{ (#1)}.}\ }}
\NewTColorBox{proposizione}{o}{nota,
  before upper={\textbf{Proposizione\IfValueT{#1}{ (#1)}.}\ }}
\NewTColorBox{proprieta}{o}{nota,
  before upper={\textbf{Proprietà\IfValueT{#1}{ (#1)}.}\ }}
\NewTColorBox{assioma}{o}{nota, boxrule=1pt,
  before upper={\textbf{Assioma\IfValueT{#1}{ (#1)}.}\ }}

% --- Ambienti semplici (senza riquadro) ---
% Il titolo sta in cima al blocco, anche se il contenuto inizia con un riquadro o
% con codice; e non resta mai da solo in fondo a una pagina.
% Uso: \begin{esempio}[Titolo facoltativo] ... \end{esempio}
\newcommand{\newrunin}[2]{%
  \NewDocumentEnvironment{#1}{o}{%
    \par\Needspace{4\baselineskip}\addvspace{\medskipamount}\noindent
    \textbf{#2}\IfValueT{##1}{ (##1)}\textbf{.}\ \ignorespaces}%
  {\par\addvspace{\medskipamount}}}
\newrunin{esempio}{Esempio}
\newrunin{esercizio}{Esercizio}
\newrunin{osservazione}{Osservazione}

% V e F nelle tavole di verità
\newcommand{\V}{\textsf{V}}
\newcommand{\F}{\textsf{F}}

% --- Codice C ---
% Uso: \begin{codice} ... \end{codice}      codice C numerato
%      \begin{comando} ... \end{comando}    comandi da terminale
%      \begin{risultato} ... \end{risultato} output di un programma
%      \lstinline|printf("ciao");|           codice dentro al testo
\lstdefinestyle{c}{
  language=C,
  basicstyle=\ttfamily\small,
  keywordstyle=\color{blue}\bfseries,
  commentstyle=\color{gray}\itshape,
  stringstyle=\color{red!70!black},
  numbers=left,
  numberstyle=\tiny\color{gray},
  numbersep=8pt,
  columns=flexible,
  keepspaces=true,
  breaklines=true,
  showstringspaces=false,
  tabsize=4,
  morekeywords={include, define},
  % lettere accentate nei commenti
  literate={à}{{\`a}}1 {è}{{\`e}}1 {é}{{\'e}}1 {ì}{{\`i}}1
           {ò}{{\`o}}1 {ù}{{\`u}}1 {È}{{\`E}}1 {À}{{\`A}}1
}
\lstset{style=c}
\newtcblisting{codice}{listing only, nota, left=6mm,
  listing options={style=c}}
\newtcblisting{comando}{listing only, nota,
  listing options={basicstyle=\ttfamily\small, numbers=none, language={},
    columns=flexible, keepspaces=true, breaklines=true}}
\newtcblisting{risultato}{listing only, enhanced jigsaw, breakable,
  colback=black!5, colframe=black, boxrule=0.5pt, sharp corners,
  listing options={basicstyle=\ttfamily\small, numbers=none, language={},
    columns=flexible, keepspaces=true, breaklines=true}}

% --- Diagrammi di flusso (simboli standard) ---
\tikzset{
  terminale/.style = {ellipse, draw, minimum width=2.5cm, minimum height=0.9cm},
  io/.style        = {trapezium, trapezium left angle=70, trapezium right angle=110,
                      draw, minimum width=2.5cm, minimum height=0.9cm},
  processo/.style  = {rectangle, draw, minimum width=2.5cm, minimum height=0.9cm},
  decisione/.style = {diamond, draw, aspect=2, minimum width=2.5cm},
  freccia/.style   = {thick, -{Stealth}}
}

\title{Insiemi}
\author{Marco Cerbella}
\date{Lezione 1 -- 21 settembre 2026}

\begin{document}
\maketitle

\section{Insiemi ed elementi}

\begin{definizione}
Un \emph{insieme} è una collezione (o famiglia) di oggetti di cui possiamo
dire con precisione se un oggetto appartiene o no alla famiglia. Un oggetto
di un insieme si chiama \emph{elemento}.
\end{definizione}

Gli insiemi si indicano di solito con lettere maiuscole ($A, B, X, S, V,
\dots$), gli elementi con lettere minuscole ($x, y, v, k, \dots$).
\begin{itemize}
    \item $x \in S$: l'elemento $x$ appartiene all'insieme $S$;
    \item $y \notin S$: $y$ non appartiene a $S$.
\end{itemize}

Un insieme si può dare in due modi.
\begin{enumerate}
    \item \textbf{Per elencazione}: si scrivono gli elementi tra graffe.
          Ad esempio $V = \{1, 3, 5, 88\}$ definisce l'insieme $V$ per
          elencazione.
    \item \textbf{Per descrizione} (o per caratteristica): ad esempio
          $X = \{\text{le lettere dell'alfabeto inglese}\}$.
\end{enumerate}

Quando la descrizione usa una formula aperta si scrive
\[
W = \{x \mid x^2 - 3x + 2 = 0\}, \qquad
A = \{x \in \mathbb{R} \mid P(x)\}, \qquad
B = \{x \in \mathbb{R} \mid x > 3\} = \{x : x > 3\},
\]
dove ``$\mid$'' (oppure ``$:$'') si legge ``tale che''.

Un insieme con un solo elemento, come $U = \{4\}$ (per cui $4 \in U$ e
$3 \notin U$), si chiama \emph{singoletto} (in inglese \emph{singleton}).

L'\emph{insieme vuoto}, indicato con $\emptyset = \{\ \}$, è l'insieme senza
elementi: per ogni $x$ vale $x \notin \emptyset$ (cioè
$\nexists\, x : x \in \emptyset$).

\section{Sottoinsiemi e uguaglianza}

\begin{definizione}
Siano $A, B$ insiemi. $B$ è un \emph{sottoinsieme} di $A$ se
\[
\forall\, x \in B \implies x \in A.
\]
Si scrive $B \subset A$ oppure $B \subseteq A$; equivalentemente
$A \supseteq B$ (o $A \supset B$).
\end{definizione}

\begin{definizione}
Due insiemi $S, T$ sono \emph{uguali}, $S = T$, se hanno gli stessi
elementi, cioè
\[
x \in S \iff x \in T.
\]
\end{definizione}

\begin{osservazione}
$S = T$ se e solo se $S \subseteq T$ e $T \subseteq S$.
\end{osservazione}

Dato un insieme $A$ valgono sempre $\emptyset \subseteq A$ e
$A \subseteq A$.

\subsection*{Insiemi numerici}

\begin{center}
\begin{tabular}{ll}
$\mathbb{N} = \{0, 1, 2, 3, \dots\}$ & numeri naturali \\
$\mathbb{Z} = \{0, \pm 1, \pm 2, \pm 3, \dots\}
              = \{\dots, -3, -2, -1, 0, 1, 2, \dots\}$ & numeri interi \\
$\mathbb{Q} = \left\{ \dfrac{a}{b} \;\middle|\; a, b \in \mathbb{Z},\
              b \neq 0 \right\}$ & numeri razionali \\
$\mathbb{R}$ & numeri reali \\
$\mathbb{C}$ & numeri complessi
\end{tabular}
\end{center}

Valgono le inclusioni
\[
\mathbb{N} \subseteq \mathbb{Z} \subseteq \mathbb{Q} \subseteq
\mathbb{R} \subseteq \mathbb{C}.
\]
Con $i$ si indica l'unità immaginaria, $i^2 = -1$. Ad esempio:
\begin{center}
\begin{tabular}{llll}
$-2 \in \mathbb{Z}$, & $\frac12 \in \mathbb{Q}$, &
$\sqrt{2}, \pi \in \mathbb{R}$, & $i \in \mathbb{C}$, \\
$-2 \notin \mathbb{N}$, & $\frac12 \notin \mathbb{Z}$, &
$\sqrt{2}, \pi \notin \mathbb{Q}$, & $i \notin \mathbb{R}$.
\end{tabular}
\end{center}

\section{Insieme delle parti}

\begin{definizione}
Sia $A$ un insieme. L'\emph{insieme delle parti} di $A$ è l'insieme di
tutti i sottoinsiemi di $A$:
\[
\mathcal{P}(A) = \{B \mid B \subseteq A\}.
\]
\end{definizione}

\begin{osservazione}
Per ogni insieme $A$: $\emptyset \subseteq A$ e $\emptyset \in
\mathcal{P}(A)$; $A \subseteq A$ e $A \in \mathcal{P}(A)$.
\end{osservazione}

\begin{esempio}
Sia $A = \{1, 2, 3\}$. Allora
\[
\mathcal{P}(A) = \big\{\emptyset,\ \{1\},\ \{2\},\ \{3\},\
\{1,3\},\ \{2,3\},\ \{1,2\},\ A\big\}.
\]
Bisogna distinguere bene appartenenza e inclusione:
\begin{itemize}
    \item $1 \in A$, \quad $\{1\} \subseteq A$, \quad
          $\{1, 3\} \subseteq A$;
    \item $\{1\} \in A$: \textbf{no}; \quad $\{1\} \in \mathcal{P}(A)$:
          sì;
    \item $A \subseteq \mathcal{P}(A)$: \textbf{no}; \quad
          $A \in \mathcal{P}(A)$: sì.
\end{itemize}
\end{esempio}

\begin{esempio}
$\mathcal{P}(\emptyset) = \{\emptyset\}$ e
\[
\mathcal{P}(\{\emptyset\}) = \mathcal{P}(\mathcal{P}(\emptyset))
= \{\emptyset, \{\emptyset\}\}.
\]
\end{esempio}

\section{Operazioni tra insiemi}

Siano $A, B$ insiemi.

\begin{definizione}
\begin{itemize}
    \item \textbf{Unione}: $A \cup B = \{x \mid x \in A \lor x \in B\}$.
    \item \textbf{Intersezione}: $A \cap B = \{x \mid x \in A \land
          x \in B\}$.
    \item \textbf{Differenza} tra $A$ e $B$:
          $A \setminus B = \{x \mid x \in A \text{ e } x \notin B\}$.
    \item \textbf{Differenza simmetrica} tra $A$ e $B$:
          $A \,\triangle\, B = (A \setminus B) \cup (B \setminus A)$.
\end{itemize}
$A$ e $B$ sono \emph{disgiunti} se $A \cap B = \emptyset$.
\end{definizione}

Esercizio: $A \,\triangle\, B = (A \cup B) \setminus (A \cap B)$.

\begin{definizione}
Se $A \subseteq U$, il \emph{complementare} di $A$ in $U$ è l'insieme
\[
A^c = U \setminus A.
\]
\end{definizione}

\begin{esempio}
Siano $A = \{1, 2, 3, 4\} = \{4, 1, 3, 2\}$ (l'ordine non conta) e
$B = \{1, 3, 6, 9, 12\}$. Allora
\begin{align*}
A \cup B &= \{1, 2, 3, 4, 6, 9, 12\}, & A \cap B &= \{1, 3\}, \\
A \setminus B &= \{2, 4\}, & B \setminus A &= \{6, 9, 12\}, \\
A \,\triangle\, B &= \{2, 4, 6, 9, 12\}.
\end{align*}
\end{esempio}

\section{Prodotto cartesiano}

\begin{definizione}
Siano $A, B$ insiemi. Il \emph{prodotto cartesiano} di $A$ e $B$ è
\[
A \times B = \{(a, b) \mid a \in A,\ b \in B\},
\]
l'insieme delle \emph{coppie ordinate}. Vale $(a, b) \neq (b, a)$ in
generale, e
\[
(x, y) = (x', y') \iff x = x' \text{ e } y = y'.
\]
\end{definizione}

\begin{esempio}
Siano $A = \{1, 2, 3\}$ e $B = \{v, w\}$. Allora
\begin{align*}
A \times B &= \{(1,v), (1,w), (2,v), (2,w), (3,v), (3,w)\}, \\
B \times A &= \{(v,1), (w,1), (v,2), (w,2), (v,3), (w,3)\},
\end{align*}
quindi $A \times B \neq B \times A$. Inoltre
\[
A \times A = A^2 = \{(1,1), (1,2), (1,3), (2,1), (2,2), (2,3), (3,1),
(3,2), (3,3)\}.
\]
\end{esempio}

Analogamente
\[
A \times B \times C = \{(a, b, c) \mid a \in A,\ b \in B,\ c \in C\},
\qquad A \times A \times A = A^3.
\]

\end{document}
